%==========================================================================
% Exemplo de artigo cientifico UNESC
%
% v0.0 versão em 20/09/2023 por Marlon Oliveira
% v1.0 versão em 25/06/2025 por Marlon Oliveira


%==========================================================================

%==========================================================================
% Pacotes utilizados pela UNESC - não delete as linhas abaixo
%==========================================================================
\documentclass{UNESC}
\usepackage[T1]{fontenc}
\usepackage[utf8]{inputenc}%\usepackage[latin1]{inputenc}
\renewcommand{\rmdefault}{phv} % Arial
\renewcommand{\sfdefault}{phv} % Arial
\usepackage[brazilian]{babel}
\usepackage{graphicx}
%\usepackage{subfigure}
\usepackage{hyperref}
\usepackage{quoting}
\usepackage[labelsep=none]{caption}
\DeclareCaptionLabelFormat{figuracomtraco}{#1~#2 -}
\captionsetup[figure]{labelformat=figuracomtraco, labelsep=space}
\captionsetup[table]{labelformat=simple, labelsep=space, textformat=simple}
\DeclareCaptionLabelFormat{tabelacomtraco}{#1~#2 -}
\captionsetup[table]{labelformat=tabelacomtraco}
\captionsetup{labelsep=none}
%\renewcommand{\thefigure}{\arabic{figure} - }
\usepackage{enumitem}
\usepackage{float}
\usepackage{comment}
\usepackage[hang,flushmargin]{footmisc}
\usepackage{array}

\hypersetup{
  setpagesize  = false,
  colorlinks   = true,    % Colours links instead of ugly boxes
  urlcolor     = blue,    % Colour for external hyperlinks
  linkcolor    = black,   % Colour of internal links
  citecolor    = black    % Colour of citations
}
\urlstyle{same}
\usepackage{setspace}

\usepackage[alf,recuo=0.0cm,abnt-emphasize=bf]{abntex2cite} 

%\captionsetup[figure]{justification=raggedright}


%==========================================================================
% Dados do artigo: título, autores e filiação
%==========================================================================
\title{TÍTULO DO ARTIGO CIENTÍFICO: SUBTÍTULO (SE HOUVER)}
\author
{
   Nome Completo do Autor\footnote{Curso de Ciência da Computação, Universidade do Extremo Sul Catarinense (Unesc), Criciúma - Santa Catarina - Brasil. email@unesc.net},
   %Nome Completo do corientador(a)\footnote{Corientador(a), Curso de Ciência da Computação, Universidade do Extremo Sul Catarinense (Unesc), Criciúma - Santa Catarina - Brasil. email@unesc.net}
   Nome Completo do Orientador\footnote{Orientador(a), Curso de Ciência da Computação, Universidade do Extremo Sul Catarinense (Unesc), Criciúma - Santa Catarina - Brasil. email@unesc.net}
}
%==========================================================================
% Início do texto do artigo
%==========================================================================
\begin{document}

\onehalfspacing

\maketitle

\noindent
\textbf{Resumo:} Deve ter uma sequência de frases concisas e objetivas em parágrafo único, sem enumeração de tópicos, de 100 a 250 palavras, seguido, logo abaixo, das palavras representativas do conteúdo do trabalho, isto é, palavras-chave e/ou descritores, conforme a \citeonline{ABNT2021}.

\vspace{0.5em}
\noindent
\textbf{Palavras-chave:} Elemento obrigatório. Devem ser separadas entre si por ponto e vírgula e finalizadas por ponto, com as iniciais em letra minúscula, com exceção dos substantivos próprios e nomes científicos. 

\vspace{1.0em}
\noindent
\textbf{ABSTRACT:} Tradução do resumo em língua vernácula para outro idioma de propagação internacional (em inglês ABSTRACT, em francês RESUMÉ, em espanhol RESUMEN).

\vspace{0.5em}
\noindent
\textbf{Keywords:} keyword 1; keyword 2; keyword 3.


\section{INTRODUÇÃO}

A introdução apresenta o assunto, delimita o tema e analisa a problemática a ser investigada, com as questões norteadoras e as hipóteses formuladas. Nela devem constar também os objetivos da pesquisa, a justificativa da escolha do tema e a metodologia utilizada.

Consiste no que se deseja atingir com o desenvolvimento da pesquisa, deve evidenciar de forma genérica o propósito mais amplo da realização do trabalho. O objetivo da pesquisa tanto o geral quanto os específicos devem ser bem definidos, já que se refere ao que será alcançado com a finalização do trabalho. Objetivos geral inicia por verbo no infinitivo.



Os objetivos específicos são o detalhamento do objetivo geral, devendo refletir o que se espera do trabalho em relação aos aspectos computacionais e ao domínio de aplicação. Estes objetivos devem ser passíveis de realização, considerando-se o tempo disponível para realização da pesquisa. Objetivos específicos iniciam por verbo no infinitivo.
Deve-se escrever sobre os motivos que determinaram a escolha do tema da pesquisa, destacando-se a relevância computacional, bem como do domínio de aplicação. 
Para facilitar o processo de redação da justificativa, deve-se perguntar: \\
-	O tema é relevante? Por quê? 

-	Quais os pontos positivos da abordagem a ser pesquisada? 

-	Que vantagens/benefícios esta pesquisa pode proporcionar? 

Exemplo de citação de referência: \cite{gandini2002} ou \citeonline{gandini2002}
Citação de site: \cite{unep}
A seguir, um exemplo com citação direta longa. \\
Segundo \citeonline[p. 57]{oliveira2017irish}: \\


\begin{quoting}[rightmargin=0cm,leftmargin=4cm]
\begin{singlespace}
{\footnotesize \noindent
A performance se estrutura, portanto, numa linguagem “cênico-teatral” e é apresentada na forma de um mixed-media onde a tonicidade maior pode dar-se um uma linguagem ou outra, dependendo da origem do artista (mais plástica no Fluxus, mais teatral em Disappearances).} 
\end{singlespace}
\end{quoting}

Caso tenha fórmulas ou equações, para facilitar a leitura, devem ser destacadas no texto e, se necessário, numeradas com algarismos arábicos entre parênteses, alinhados à direita. Na sequência normal do texto, é permitido o uso de uma entrelinha maior que comporte seus elementos (expoentes, índices, entre outros). \\
$x2 +  y2 = z2$ \\
$(x2 +  y2)/5 = n$ 

\begin{equation}
acuracia = \frac{numero \ de \ acertos}{numero de tentativas}
\label{eq:pitagoras}
\end{equation}

Exemplo de como inserir Figura \ref{fig:fig1} e Tabela \ref{tab:table1}. 

\begin{figure}[H]
\centering  
\caption{Acervo da Biblioteca Central da UNESC}
\vspace{-0.2cm} 
\label{fig:fig1}
\includegraphics[scale=0.3]{Figs/Fig1.jpeg} 
\flushleft
\vspace{-0.2cm} 
\hspace{0.9cm} \footnotesize Fonte: Elaborado pelo autor.
\end{figure}

\begin{table}[h!]
\begin{center}
\caption{Pessoas residentes em domicílios particulares}
\vspace{-0.2cm} 
    \begin{center}
    \label{tab:table1}
    %\resizebox{\textwidth}{!}{%
    \begin{tabular}{l l l l} 
    \hline
    Situação do domicílio & Total & Mulheres & Homens \\
    \hline
    Total   & 117.960.301 & 59.595.332 & 58.364.969 \\
    Urbana  & 79.972.931  & 41.115.439 & 38.857.492 \\
    Rural   & 37.987.370  & 18.479.893 & 19.507.477 \\
    \hline
    \end{tabular} 
  \end{center}
\end{center}  
\vspace{-0.2cm}
\hspace{1cm}\footnotesize Fonte: \citeonline{gandini2002}.
\end{table}

\section{MATERIAIS E MÉTODOS}

\section{RESULTADOS E DISCUSSÃO}


\section{Conclusão}
%Parte final do artigo, onde são apresentadas as conclusões correspondentes aos objetivos e as hipóteses. Deve ser breve, devendo apresentar recomendações e sugestões para trabalhos futuros.

\bibliographystyle{abntex2-unesc}
\renewcommand{\refname}{REFERÊNCIAS}
\bibliography{referencias} % arquivo referencias.bib que contém as referências bibliográficas utilizadas no artigo

%\include{apendice_a}
%\include{anexo_a}

\end{document} 